% ============================================================
%  Distance Project — Research Storyboard / Results Summary
%  Compile: pdflatex storyboard.tex  (twice for TOC/refs)
% ============================================================
\documentclass[11pt, letterpaper]{article}

% ── Layout ───────────────────────────────────────────────────
\usepackage[margin=1in]{geometry}
\usepackage{setspace}
\onehalfspacing

% ── Fonts & encoding ─────────────────────────────────────────
\usepackage[T1]{fontenc}
\usepackage{microtype}
\usepackage{lmodern}

% ── Colour ───────────────────────────────────────────────────
\usepackage[dvipsnames]{xcolor}
\definecolor{colRed}{HTML}{C0392B}
\definecolor{colBlue}{HTML}{2E86AB}
\definecolor{colGrey}{HTML}{555555}
\definecolor{colLight}{HTML}{F4F6F7}
\definecolor{colRule}{HTML}{CCCCCC}

% ── Tables ───────────────────────────────────────────────────
\usepackage{booktabs}
\usepackage{tabularx}
\usepackage{array}
\usepackage{multirow}

% ── Section styling ──────────────────────────────────────────
\usepackage{titlesec}
\titleformat{\section}
  {\large\bfseries\color{black}}
  {\thesection.}{0.6em}{}
  [\vspace{-0.3em}\color{colRule}\rule{\linewidth}{0.6pt}]
\titleformat{\subsection}
  {\normalsize\bfseries\color{colGrey}}
  {\thesubsection}{0.6em}{}

% ── Callout boxes ────────────────────────────────────────────
\usepackage{tcolorbox}
\tcbuselibrary{skins, breakable}

\tcbset{
  finding/.style={
    enhanced, breakable,
    colback=colLight, colframe=colRed,
    fonttitle=\bfseries\small,
    title={Key Finding},
    left=6pt, right=6pt, top=4pt, bottom=4pt,
    boxrule=1pt, arc=3pt
  },
  note/.style={
    enhanced, breakable,
    colback=colLight, colframe=colBlue,
    fonttitle=\bfseries\small,
    left=6pt, right=6pt, top=4pt, bottom=4pt,
    boxrule=0.8pt, arc=3pt
  }
}

% ── Misc ─────────────────────────────────────────────────────
\usepackage{parskip}
\usepackage{graphicx}
\usepackage{caption}
\usepackage{hyperref}
\hypersetup{colorlinks=true, linkcolor=colBlue, urlcolor=colBlue}

% ── Helpers ──────────────────────────────────────────────────
\newcommand{\red}[1]{{\color{colRed}\textbf{#1}}}
\newcommand{\wht}[1]{{\color{colBlue}\textbf{#1}}}
\newcommand{\sig}[1]{{\color{colRed}#1}}
\newcommand{\ns}[1]{{\color{colGrey}#1}}
\newcommand{\pval}[1]{\textit{p}\,$=$\,#1}
\newcommand{\figbox}[1]{%
  \begin{tcolorbox}[note, title=Figure placeholder]
    \small\color{colGrey} #1
  \end{tcolorbox}}

% ============================================================
\begin{document}
% ============================================================

% ── Title ────────────────────────────────────────────────────
\begin{center}
  {\LARGE\bfseries Artificial Light at Night and Bat Activity\\[0.3em]
   Along a Distance Gradient}\\[0.8em]
  {\large Grand Teton National Park, 2022}\\[0.5em]
  {\normalsize Nigel K.\ Anderson \quad $\cdot$ \quad
   American Museum of Natural History}\\[0.3em]
  {\small\color{colGrey}\today}
\end{center}

\vspace{0.5em}
\noindent\color{colRule}\rule{\linewidth}{1pt}
\vspace{0.5em}

% ── Abstract / Overview ──────────────────────────────────────
\noindent\textbf{Overview.}
Most research on artificial light at night (ALAN) compares light versus dark.
This study asked a more nuanced question: does the \textit{colour} of light
matter for bats, and does that depend on \textit{how bright} it is or
\textit{how far away} you are from the source?
We monitored acoustic bat activity at 14 sites along a distance gradient
in Grand Teton National Park, crossing \red{red} versus \wht{white} light
at five intensity levels (10, 30, 50, 70, 100\,\%).
The results reveal a strong colour $\times$ intensity interaction,
a near-complete collapse of bat activity close to the light source,
and a clear separation between the drivers of total activity
(colour $\times$ intensity) and community composition (distance, habitat).

\tableofcontents
\vspace{1em}

% ============================================================
\section{Study Design and Statistical Approach}
% ============================================================

Fourteen acoustic monitoring stations were deployed across a distance
gradient ranging from $<$0.52\,km (\textit{Close}) to $>$2.1\,km
(\textit{Far}) from a central experimental light source.
Nightly detection data were collected across the 2022 survey season.
Eight focal bat species were retained for analysis:
\textit{Epfu}, \textit{Laci}, \textit{Lano}, \textit{Myev},
\textit{Mylu}, \textit{Myvo}, \textit{Myyu}, and \textit{Myci}.

The primary statistical model was a zero-inflated negative binomial GLMM
(glmmTMB) with a random intercept for site:

\begin{center}
\texttt{detections $\sim$ color $\times$ intensity + moon\_phase +
\%nonforest + sky\_brightness$_\text{dark}$ + jd$_c$ + jd$_c^2$ +
dist\_cat + (1\,|\,site)}
\end{center}

Julian day was mean-centred (\texttt{jd}$_c$) to avoid collinearity
in the quadratic term.
Post-hoc contrasts (emmeans, Holm-corrected) were computed on the
response scale (detections per night).
Community-level analyses used NMDS, PERMANOVA, distance-based redundancy
analysis (dbRDA), and multivariate GLMs (mvabund).

% ============================================================
\section{Overall Model Fit}
% ============================================================

The model converged with AIC\,$=$\,24{,}276 on 3,457 observations across
14 sites ($n_\text{groups}$\,=\,14, $\sigma_\text{site}$\,=\,0.55).
The negative binomial dispersion parameter was 0.603, consistent with
overdispersion in count data. Table~\ref{tab:model} summarises all
fixed-effect estimates.

\begin{table}[h]
\centering
\caption{Fixed-effect estimates from the zero-inflated negative binomial
GLMM. Reference levels: \texttt{color} = Red, \texttt{intensity} = 10\,\%,
\texttt{dist\_cat} = Close.}
\label{tab:model}
\smallskip
\begin{tabular}{lrrrr}
\toprule
\textbf{Term} & \textbf{Estimate} & \textbf{SE} & \textbf{\textit{z}} & \textbf{\textit{p}} \\
\midrule
(Intercept)          &  1.286 & 0.608 &  2.12 & 0.034\,* \\
colorW               &  0.011 & 0.092 &  0.12 & \ns{0.908} \\
intensity30          &  0.596 & 0.137 &  4.37 & \sig{$<$0.001\,***} \\
intensity50          &  0.354 & 0.143 &  2.48 & \sig{0.013\,*} \\
intensity70          &  0.354 & 0.093 &  3.79 & \sig{$<$0.001\,***} \\
intensity100         & $-$0.000 & 0.116 & $-$0.00 & \ns{0.998} \\
moon phase           &  0.215 & 0.116 &  1.86 & \ns{0.063\,.} \\
\%\,nonforest        &  2.130 & 1.133 &  1.88 & \ns{0.060\,.} \\
sky brightness (dark)& $-$0.026 & 0.034 & $-$0.78 & \ns{0.438} \\
jd$_c$               & $-$0.007 & 0.002 & $-$4.25 & \sig{$<$0.001\,***} \\
jd$_c^2$             & $-$0.000 & 0.000 & $-$1.70 & \ns{0.088\,.} \\
dist\_cat: Far       &  0.990 & 0.482 &  2.06 & \sig{0.040\,*} \\
dist\_cat: Further   &  1.208 & 0.599 &  2.02 & \sig{0.044\,*} \\
dist\_cat: Medium    &  1.196 & 0.676 &  1.77 & \ns{0.077\,.} \\
\midrule
colorW:intensity30   & $-$0.728 & 0.171 & $-$4.26 & \sig{$<$0.001\,***} \\
colorW:intensity50   & $-$0.497 & 0.146 & $-$3.42 & \sig{$<$0.001\,***} \\
colorW:intensity70   & $-$0.123 & 0.136 & $-$0.90 & \ns{0.368} \\
colorW:intensity100  &  0.089 & 0.134 &  0.67 & \ns{0.505} \\
\bottomrule
\end{tabular}
\end{table}

% ============================================================
\section{Colour and Intensity Effects}
% ============================================================

\subsection{Main effect of colour}

Averaged across all intensity levels and distance categories, \red{red}
and \wht{white} light produced statistically identical bat activity
($\beta_\text{white}$\,=\,0.011, \pval{0.91}).
The headline result is not about colour per se --- it emerges only in
interaction with intensity.

\subsection{The colour $\times$ intensity interaction}

\begin{tcolorbox}[finding]
  At 30\,\% intensity, red light produced \textbf{2.05$\times$ more}
  bat detections than white light (\pval{$<$0.0001}).
  At 50\,\%, the ratio was 1.63$\times$ (\pval{$<$0.0001}).
  At 70\,\% and 100\,\%, colours were statistically indistinguishable.
  \textbf{White light produced the same activity level regardless of intensity.}
\end{tcolorbox}

Table~\ref{tab:emm} presents the emmeans-estimated marginal means and
pairwise colour contrasts on the response scale.

\begin{table}[h]
\centering
\caption{Estimated marginal means (detections/night) by colour and
intensity, with Holm-corrected pairwise contrasts.
Averaged over distance categories.}
\label{tab:emm}
\smallskip
\begin{tabular}{lccccc}
\toprule
\textbf{Intensity} & \red{Red} & \wht{White}
                   & \textbf{95\,\% CI (Red)} & \textbf{Ratio R/W}
                   & \textbf{\textit{p}} \\
\midrule
10\,\% & 13.1 & 13.3 & [9.4,\;18.4]  & 0.99 & \ns{0.908} \\
30\,\% & \textbf{23.8} & \textbf{11.6} & [16.3,\;34.7] & \textbf{2.05} & \sig{$<$0.001} \\
50\,\% & \textbf{18.7} & \textbf{11.5} & [12.5,\;27.8] & \textbf{1.63} & \sig{$<$0.001} \\
70\,\% & 18.7 & 16.7 & [13.3,\;26.2] & 1.12 & \ns{0.257} \\
100\,\%& 13.1 & 14.5 & [9.2,\;18.7]  & 0.91 & \ns{0.317} \\
\bottomrule
\end{tabular}
\end{table}

\figbox{Fig.\ 3 --- Connected dot plot: mean detections by colour at each
intensity level with 95\,\% CI.
Red line peaks at 30\,\% and converges with white at high intensity.}

\subsection{Intensity within each colour}

Within \red{red} light, 30\,\% intensity was the peak activity condition
(23.8 det/night), significantly exceeding near-dark (10\,\%) by 82\,\%
(\pval{0.0001}) and 100\,\% intensity by 82\,\% (\pval{$<$0.001}).
The 30\,\%, 50\,\%, and 70\,\% levels did not differ significantly from
each other.

Within \wht{white} light, \textbf{no intensity level differed from any
other} (all Holm-corrected $p \geq 0.18$).
White light consistently produced $\approx$11--17 detections per night
regardless of brightness.

\begin{tcolorbox}[note]
\textbf{Interpretation.}
Moderate red light likely attracts insects, and bats follow.
White light suppresses this insect-mediated effect at any intensity.
At maximum intensity (100\,\%), both colours return to near-dark baseline ---
likely because bats abandon the area entirely rather than adjusting their
foraging behaviour.
\end{tcolorbox}

\figbox{Fig.\ 1 --- Bar chart: mean $\pm$ SE detections by colour
(averaged across intensities). Red $\approx$ White when collapsed.}

% ============================================================
\section{Distance Effects}
% ============================================================

Sites close to the light source ($<$0.52\,km) had markedly fewer bat
detections than all other distance categories
(dist\_cat: Far $\beta = 0.99$, \pval{0.040};
Further $\beta = 1.21$, \pval{0.044};
Table~\ref{tab:dist}).
The community-level analysis confirmed distance as the strongest single
predictor of community composition (dbRDA $F = 8.03$, \pval{$<$0.001}).

\begin{table}[h]
\centering
\caption{Emmeans-estimated marginal means by distance category
(averaged over colour and intensity).
Pairwise contrasts are Holm-corrected; none survive correction,
but the magnitude of the Close vs.\ all-other difference is
biologically substantial.}
\label{tab:dist}
\smallskip
\begin{tabular}{lcccc}
\toprule
\textbf{Category} & \textbf{Range (km)} & \textbf{Det/night}
                  & \textbf{95\,\% CI} & \textbf{vs.\ Close (ratio)} \\
\midrule
Close   & $<$0.52   & \red{6.5}  & [2.9,\;14.2] & --- \\
Medium  & 0.52--1.0 & 21.4 & [9.0,\;50.5] & 3.30$\times$ \\
Further & 1.0--2.1  & 21.6 & [11.7,\;40.0] & 3.34$\times$ \\
Far     & $>$2.1    & 17.4 & [8.8,\;34.2] & 2.69$\times$ \\
\bottomrule
\end{tabular}
\end{table}

\begin{tcolorbox}[finding]
Sites within 0.52\,km of the light source had $\approx$\textbf{3$\times$
fewer bat detections} than all other distances.
Medium, Further, and Far sites were indistinguishable from each other,
suggesting near-complete recovery of bat activity beyond 0.5\,km.
\end{tcolorbox}

The Holm correction is conservative given the small number of sites and
resulting low power; the main model and dbRDA both confirm the distance
effect independently. Pairwise contrasts should be interpreted alongside
the descriptive means rather than relied on exclusively.

\figbox{Fig.\ 2 --- Bar chart: mean $\pm$ SE detections by dist\_cat.\\
Fig.\ 7 --- Site-level scatter: dist\_km vs.\ mean detections with
linear trend and Spearman $\rho$.}

% ============================================================
\section{Seasonal and Environmental Covariates}
% ============================================================

Bat activity declined significantly through the survey season
(jd$_c$: $\beta = -0.007$, \pval{$<$0.001}), with a marginal
quadratic term suggesting the decline was somewhat non-linear
(\pval{0.088}).
Moon phase was marginally positive ($\beta = 0.215$, \pval{0.063}),
possibly reflecting insect availability rather than direct bat behaviour.
Neither sky brightness nor percent non-forest reached significance in
the main model (both $p > 0.05$), though both were significant
community-level predictors (see Section~\ref{sec:community}).

\figbox{Fig.\ 4 --- Seasonal LOESS by colour.
Red/white separation persists across the survey period.\\
Fig.\ 6 --- Moon phase LOESS by colour.
No evident colour $\times$ moon interaction.}

% ============================================================
\section{Community Composition}
\label{sec:community}
% ============================================================

The PERMANOVA explained 25.6\,\% of variance in community composition
($R^2 = 0.256$, $F = 3.75$, \pval{$<$0.001}).
The dispersion test was non-significant (\pval{0.938}), confirming that
communities are equally variable between colour treatments ---
the PERMANOVA result reflects genuine compositional shifts in location,
not differences in spread.

\begin{table}[h]
\centering
\caption{Marginal effects on bat community composition from dbRDA and
mvabund multivariate GLM. Significant predictors ($p < 0.05$)
in bold.}
\label{tab:community}
\smallskip
\begin{tabular}{lcccc}
\toprule
\textbf{Predictor} & \textbf{dbRDA $F$} & \textbf{dbRDA \textit{p}}
                   & \textbf{mvabund Dev} & \textbf{mvabund \textit{p}} \\
\midrule
colour               & 1.45 & \ns{0.196}          &   9.2 & \ns{0.394} \\
intensity            & 1.69 & \sig{0.046*}         &  83.5 & \sig{0.002**} \\
\textbf{dist\_cat}   & \textbf{8.03} & \sig{$<$0.001***} & \textbf{311.9} & \sig{$<$0.001***} \\
moon phase           & 0.86 & \ns{0.488}          &  13.2 & \ns{0.197} \\
\textbf{\%\,nonforest}& \textbf{7.53} & \sig{$<$0.001***} & \textbf{42.2} & \sig{$<$0.001***} \\
\textbf{sky brightness}& \textbf{4.20} & \sig{0.003**} & \textbf{57.4} & \sig{$<$0.001***} \\
\bottomrule
\end{tabular}
\end{table}

\begin{tcolorbox}[finding]
  \textbf{Colour does not restructure bat community composition.}
  Distance does --- by a wide margin.
  Intensity shifts composition weakly.
  Habitat openness and sky darkness independently shape who is present,
  independent of the experimental light treatment.
\end{tcolorbox}

\figbox{Fig.\ C1 --- dbRDA biplot. Distance arrows should dominate;
colour arrow should be short.\\
Fig.\ C2 --- Species activity heatmap: colour $\times$ intensity.\\
Fig.\ C3 --- Species activity heatmap: distance categories.}

% ============================================================
\section{Species-Level Responses}
% ============================================================

\subsection{Univariate community responses}

Table~\ref{tab:species} summarises species-level p-values from the
mvabund univariate tests (adjusted for multiple comparisons).

\begin{table}[h]
\centering
\caption{Univariate mvabund test results by species. Adjusted
\textit{p}-values; \sig{***}\,$<$\,0.001, \sig{**}\,$<$\,0.01,
\sig{*}\,$<$\,0.05, \ns{ns}\,$\geq$\,0.05.}
\label{tab:species}
\smallskip
\begin{tabular}{lccccc}
\toprule
\textbf{Species} & \textbf{Colour} & \textbf{Intensity}
                 & \textbf{Distance} & \textbf{\%\,nonforest}
                 & \textbf{Sky brightness} \\
\midrule
\textit{Epfu} & \ns{ns} & \sig{***} & \sig{***} & \ns{ns}    & \ns{ns} \\
\textit{Laci} & \ns{ns} & \ns{.}    & \sig{**}  & \ns{ns}    & \ns{ns} \\
\textit{Lano} & \ns{ns} & \ns{ns}   & \sig{**}  & \ns{ns}    & \ns{ns} \\
\textit{Myev} & \ns{ns} & \ns{ns}   & \sig{***} & \ns{.}     & \ns{ns} \\
\textit{Mylu} & \ns{ns} & \ns{ns}   & \sig{**}  & \sig{***}  & \ns{ns} \\
\textit{Myvo} & \ns{ns} & \ns{ns}   & \sig{***} & \ns{ns}    & \sig{***} \\
\textit{Myyu} & \ns{ns} & \ns{ns}   & \ns{ns}   & \ns{ns}    & \ns{ns} \\
\textit{Myci} & \ns{ns} & \ns{ns}   & \ns{ns}   & \ns{ns}    & \ns{ns} \\
\bottomrule
\end{tabular}
\end{table}

\textit{Epfu} is the most treatment-sensitive species, responding to
both intensity and distance.
\textit{Mylu} tracks habitat openness; \textit{Myvo} tracks sky darkness.
\textit{Myyu} and \textit{Myci} did not respond significantly to any
measured predictor.

\subsection{Per-species GLMMs}

Per-species zero-inflated negative binomial models were fit for each
focal species with sufficient data ($\geq$30 detections, $\geq$20 nights).
Results include emmeans contrasts of \red{red} vs.\ \wht{white} at each
intensity level and pairwise intensity comparisons within each colour.

\figbox{Fig.\ S1 --- Marginal means $\pm$ 95\,\% CI: colour effect ranked by species.\\
Fig.\ S2 --- White/red ratio by intensity level, per species.\\
Fig.\ S3 --- Log$_2$ ratio heatmap: White\,/\,Red across species and intensity.\\
Fig.\ 5  --- Species $\times$ colour heatmap (z-scored activity).\\
Fig.\ 8  --- Species $\times$ distance heatmap.}

% ============================================================
\section{Synthesis and Management Implications}
% ============================================================

These results support three separable conclusions.

\textbf{1. Colour matters, but only at moderate intensity.}
Red light at 30--50\,\% intensity generated substantially more bat
detections than white light at the same intensity, most likely through
insect-mediated attraction.
The suppression of this effect by white light was strong and consistent
across species.
At maximum intensity (100\,\%), both colours converged to near-dark
baseline levels, suggesting complete behavioural exclusion at that threshold.

\textbf{2. Distance is the primary driver of community structure.}
Light colour and intensity affect how many bats are detected, but do not
selectively filter species.
Distance does: sites within 0.52\,km of the light source had 3$\times$
fewer detections and a measurably different community composition.
Recovery to ambient activity levels appeared to occur between the Close
and Medium categories, suggesting a practical buffer distance of
$\approx$0.5--1.0\,km.

\textbf{3. Habitat and sky darkness independently structure the community.}
Percent non-forest and natural sky brightness were significant community-level
predictors independent of the experimental treatment.
This means the impact of a given light installation will depend strongly
on where it is placed --- open habitat near naturally dark skies will
be more severely affected than forested or already-lit areas.

\bigskip
\begin{tcolorbox}[note, title={Management Recommendations}]
\begin{enumerate}\setlength\itemsep{0.4em}
  \item Use \red{red} light only where insect-mediated bat attraction is
        acceptable or desirable; use \wht{white} where suppression of bat
        activity is needed.
  \item Keep artificial light sources \textbf{$>$0.5\,km from core bat
        foraging and commuting habitat}.
  \item Avoid 30--50\,\% intensity red light near sensitive sites (roosts,
        water, edge habitat) --- this is the highest-impact combination.
  \item Monitor \textit{Epfu} as a sentinel species: it responds to both
        intensity and proximity and may serve as an early indicator of
        community-level ALAN impacts.
  \item Site new infrastructure with reference to natural sky brightness
        and habitat openness, not light treatment alone.
\end{enumerate}
\end{tcolorbox}

% ============================================================
\section*{Species Key}
% ============================================================
\addcontentsline{toc}{section}{Species Key}

\begin{tabular}{ll}
\toprule
\textbf{Code} & \textbf{Species} \\
\midrule
\textit{Epfu} & \textit{Eptesicus fuscus} (big brown bat) \\
\textit{Laci} & \textit{Lasiurus cinereus} (hoary bat) \\
\textit{Lano} & \textit{Lasionycteris noctivagans} (silver-haired bat) \\
\textit{Myci} & \textit{Myotis ciliolabrum} (western small-footed myotis) \\
\textit{Myev} & \textit{Myotis evotis} (long-eared myotis) \\
\textit{Mylu} & \textit{Myotis lucifugus} (little brown myotis) \\
\textit{Myvo} & \textit{Myotis volans} (long-legged myotis) \\
\textit{Myyu} & \textit{Myotis yumanensis} (Yuma myotis) \\
\bottomrule
\end{tabular}

\end{document}
